\section{Related work}
\label{sec:related}

\paragraph{Run-to-run variance.} That reinforcement learning results move under reseeding is old
news in deep RL \citep{henderson2018matters} and has been documented again for language-model
reasoning, where sampling parameters, prompt formatting and hardware all shift reported numbers
\citep{hochlehnert2025sober}. Those papers establish the size of the problem by replication, which
is the only instrument available without a model of where the variance comes from. This paper
supplies such a model, and the practical difference is between paying $N$ runs and paying one run
plus a short fork.

\paragraph{Stochastic dynamics of training.} Treating SGD as a stochastic process with a stationary
distribution set by the balance of gradient noise against curvature is standard
\citep{mandt2017sgd}, and the Ornstein-Uhlenbeck argument in Section~\ref{sec:seeds} is that
argument applied to the gap between two runs. What is not established, and what
Section~\ref{sec:experiments} tests, is whether it survives RLVR: the objective is non-stationary
because the data distribution is the policy's own samples, the noise covariance depends on the
pass-rate distribution and moves as the policy improves, and the reward is binary. None of that is
in the setting the original analysis assumes, and the consequence for reproducibility -- that the
spread across seeds does not grow with the length of a run -- does not appear to have been stated
or measured.

\paragraph{Batch size and gradient noise.} The gradient noise scale predicts the batch size beyond
which data parallelism stops buying speed, and does so from statistics measurable during training
\citep{mccandlish2018empirical}; the empirical picture across tasks and optimisers is mapped in
\citet{shallue2019measuring}. Both concern a single level of sampling. RLVR samples prompts and then
samples responses within each prompt, and the resulting two-level noise scale, the split it implies,
and the fact that the response-level term is readable off the reward histogram, have no counterpart
in that literature. The variance decomposition itself is the classical two-stage sampling result
\citep{cochran1977sampling}, applied here to policy gradients.

\paragraph{Hyperparameter transfer.} $\mu$P makes learning rates transfer across width by fixing what
a step does to the network's functions rather than to its parameters
\citep{yang2021tensor, yang2022mutransfer}. This paper takes the same stance one level up: fix what
a step does to the policy's \emph{behaviour}, measured as drift in nats per token, and the recipe
transfers across batch size. We find no comparable prescription for RL post-training in the
literature; the closest work reports empirical scaling behaviour and fitted power laws
\citep{tan2025scalingbehaviors, khatri2025artofscaling, ghosh2026predictable} rather than a
quantity to hold fixed.

\paragraph{Estimators for RLVR.} RLOO \citep{ahmadian2024rloo, williams1992reinforce}, GRPO
\citep{shao2024deepseekmath}, its unnormalised correction \citep{liu2025drgrpo}, DAPO
\citep{yu2025dapo} and sequence-level variants \citep{zheng2025gspo} are usually compared by
benchmark deltas. Proposition~\ref{prop:mean} shows that, for binary rewards, the members of this
family differ only through a difficulty weight $\lambda(p, G)$, that a $p$-independent $\lambda$ is
absorbed exactly by the step size, and that standardisation is a change of objective rather than a
variance reduction -- a sharpening of the argument in \citet{liu2025drgrpo}. Recent analyses of why
normalisation helps \citep{ge2026normalization} and of what outcome-reward objectives cannot
simultaneously be \citep{ding2026impossibility} are complementary: they characterise the estimator,
while the question here is how many samples to give it and how far to step.

\paragraph{Allocating rollouts.} Several recent methods vary the number of responses per prompt
within a batch using heuristics tied to observed pass rates or gradient variance
\citep{nguyen2026allocation}, or adapt the batch size from a divergence signal
\citep{park2026batchscaling}, and one line of work argues for scaling the response count
specifically \citep{hu2025brorl}. These are allocation policies; what has been missing is the
quantity being optimised. Theorem~\ref{thm:split} supplies it, and the answer is a ratio of two measurable traces. This
also explains why heuristics keyed to pass rate alone succeed on some corpora and not others:
the pass rate determines $\tau_w$ and says nothing about $\tau_b$. The two lines of work are
largely orthogonal, since those methods distribute responses across prompts at a group size
someone else chose, while the question here is what that group size should be; the baseline we
compare against is therefore the fixed $G$ of a published recipe rather than an allocation rule.

\paragraph{Trust regions and KL control.} Constraining how far a policy moves is old
\citep{schulman2015trpo, schulman2017ppo} and its modern LLM forms clip importance ratios or
penalise divergence from a reference. The drift used here is neither: it is the KL between
\emph{consecutive} policies, used as a measurement and a control target rather than as a penalty,
and its role in this paper is to supply the unit in which a recipe transfers.

\paragraph{Asynchrony and off-policy correction.} A parallel literature studies what happens when
the sampling policy and the training policy differ, whether through pipeline staleness
\citep{song2026staleness} or numerical mismatch between rollout and training engines. That work
asks how much mismatch a run tolerates. This paper asks a question that arises before any mismatch:
given that the sampler and the trainer agree, how should the sampling budget be divided, and how
large a step should the result support.
